\documentclass[12pt,leqno]{article}

\usepackage{amsmath,amssymb,amsthm,colortbl,xcolor,nicefrac,setspace}
\usepackage[T1]{fontenc}
\usepackage[expert,seriftt,lucidasmallscale]{lucidabr}
\usepackage{mathrsfs,hyperref,color,graphicx}
\usepackage{longtable,booktabs,array}

\DeclareMathSymbol{\Dashv}{\mathrel}{symbols}{239}
\DeclareMathSymbol{\eqdef}{\mathrel}{symbols}{214}
\DeclareMathSymbol{\sentails}{\mathrel}{symbols}{240}

\DeclareSymbolFont{stmry}{U}{stmry}{m}{n}
\DeclareMathDelimiter{\llbracket}{\mathopen}{stmry}{"4A}{stmry}{"71}
\DeclareMathDelimiter{\rrbracket}{\mathclose}{stmry}{"4B}{stmry}{"79}

\setlength{\textwidth}{6.5in}
\setlength{\textheight}{9in}
\setlength{\oddsidemargin}{0mm}
\setlength{\evensidemargin}{0mm}
\setlength{\topmargin}{0mm}
\setlength{\headheight}{0mm}
\setlength{\headsep}{0mm}
\setlength{\topskip}{0mm}
\setlength{\hoffset}{0in}
\setlength{\voffset}{0in}

\newcommand{\Notz}{\ensuremath{\mathord{\thicksim}}}
\newcommand{\AND}{\ensuremath{\mathbin{\&}}}
\newcommand{\OR}{\ensuremath{\vee}}
\newcommand{\IFF}{\ensuremath{\leftrightarrow}}
\newcommand{\IC}{\ensuremath{\rightarrow}}
\newcommand{\IF}{\ensuremath{\supset}}
\newcommand{\NOT}{\ensuremath{\text{$\thicksim$}}}
\newcommand{\given}{\ensuremath{\mathbin{|}}}
\newcommand{\entails}{\ensuremath{\vDash}}
\newcommand{\nentails}{\ensuremath{\nvDash}}
\newcommand{\dentails}{\ensuremath{\vdash}}
\newcommand{\ndentails}{\ensuremath{\nvdash}}
\newcommand{\sem}[1]{\ensuremath{\llbracket #1 \rrbracket}}

\theoremstyle{plain}
\newtheorem{theorem}{Theorem}
\newtheorem{proposition}{Proposition}
\newtheorem{corollary}{Corollary}
\theoremstyle{definition}
\newtheorem{definition}{Definition}
\newtheorem{remark}{Remark}

\onehalfspacing

\newcommand{\isa}[1]{\texttt{\detokenize{#1}}}
\newcommand{\funp}{\ensuremath{\mathsf{fun}'\,}}
\newcommand{\Pure}{\ensuremath{\mathsf{Pure}}}
\newcommand{\Fun}{\ensuremath{\mathsf{Fun}}}
\newcommand{\PP}{\ensuremath{\mathsf{PP}}}
\newcommand{\Ltwo}{\ensuremath{\mathsf{L2}}}
\newcommand{\QSS}{\ensuremath{\mathsf{QSS}}}
\newcommand{\TU}{\ensuremath{\mathsf{TU}}}
\newcommand{\WI}{\ensuremath{\mathsf{WI}}}
\newcommand{\Inv}{\ensuremath{\mathsf{Inv}}}
\newcommand{\RS}{\ensuremath{\mathsf{RS}}}
\newcommand{\PC}{\ensuremath{\mathsf{PC}}}
\newcommand{\Range}{\ensuremath{\mathsf{Range}}}
\newcommand{\N}{\ensuremath{\mathbb{N}}}
\newcommand{\Pow}{\ensuremath{\mathcal{P}}}

\hypersetup{
  colorlinks=true,
  linkcolor=blue!50!black,
  urlcolor=blue!50!black,
  pdftitle={Verification and progress on Purity of Pure},
  pdfsubject={Goodman's PP consistency question}
}

\begin{document}

\begin{center}
{\Large\bfseries Verification and Progress on Purity of Pure\par}
\vspace{4pt}
{\large A machine-checked report on Goodman's project notes\par}
\vspace{8pt}
{\normalsize 27 July 2026; updated 28 July 2026}
\end{center}

\bigskip

\noindent\textbf{Purpose.}
This report records what we have verified in Jeremy Goodman's July 2026
project notes on the consistency of Purity of Pure with Bacon's logical
combinatorialism.  It has two aims.  First, it separates the claims that have
been verified exactly from claims that require qualifications, repairs, or
further formulation.  Second, it states the new theorems proved during the
formalization that genuinely narrow the open consistency problem.

\medskip

\noindent\textbf{Bottom line.}
Goodman's central consistency question remains open.  Isabelle nevertheless
settles the determinate subsidiary claims in the notes: each is proved,
refuted and repaired, or assigned the explicit qualification under which it
is true.  In particular, the formalization identifies and repairs a modal gap
in T3 and a false countability argument in M5; verifies the four T6
contradiction routes and the T9 cardinal dichotomy from their explicit
premises; completes the rebuilt M5 model containing the displayed exotic
operator; and isolates the assumptions still needed for several of the
remaining model-theoretic claims.  It also proves that Recombination together with
zeroary Exhaustion and unique proposition-level fundamentality supplies the
\(\funp\) witness used in T6.  Hence T6 has the following exact semantic
consequence: any model of the repaired background theory plus \PP\ must
falsify global \Ltwo\ or global \TU.  That consequence sharply constrains a
positive answer, but it does not decide whether such a model exists.  We have
also settled Goodman's proposed calibration of \Ltwo\ in Bacon's appendix
model: a closed logical operator \(Z\) refutes \Ltwo\ when the pure unary
operators are the denotations of closed terms containing only logical
vocabulary.  This is a substantive result about
the notes, although Bacon's model still fails \PP\ and therefore does not
answer the consistency question.

\section{Formal setting and verification standard}

The formal development represents Bacon's typed language and proof system
directly in Isabelle/HOL.  It includes Bacon's system \(H\), classical and Boolean
identity principles, the Rule of Equivalence, vector Equivalence, the
definition of necessity by identity with truth, and the additional
purity/fundamentality vocabulary used in Goodman's notes.  The development
distinguishes object-language derivability from validity in a particular
model.  This matters throughout: a semantic argument valid in Bacon's
appendix model is not thereby a derivation in the theory, and an Isabelle
theorem with explicit premises does not show that those premises hold in the
intended model.

The complete audit checks the theorem objects supporting the claims summarized
below.  Fresh builds pass.  None of these proofs contains an admitted lemma,
an untrusted proof step, or an undischarged assumption beyond the assumptions
stated in the theorem itself.  The type restriction carried by the general
cardinal theorems in T9 is part of their formally checked statement, not a
further mathematical premise.

\section{Background definitions and two important distinctions}

At proposition type, Goodman's definable surrogate for fundamentality is
\[
 \funp(p)
 \quad\eqdef\quad
 \forall X\,Y\bigl((\Pure(X)\AND\Pure(Y)\AND Xp=Yp)\IC X=Y\bigr).
\]
Let \(G\) be the pure reversible unary operators, and let
\[
 X\approx Y
 \quad\eqdef\quad
 \exists Z\in G\;(X=Y\circ Z).
\]
Weak \Ltwo\ says that if \(X\) and \(Y\) are pure, \(p\) and \(q\) are
\(\funp\), and \(Xp=Yq\), then \(X\approx Y\).

Two distinctions uncovered by the formalization are central.

\paragraph{Recombination versus Exhaustion.}
Unary Recombination proves a pointwise modal identity principle.  It does not,
by itself, permit necessitating the universal closure required by the usual
\QSS\ derivation.  Zeroary Exhaustion supplies exactly that missing step.  With
unique proposition-level fundamentality, the repaired package derives the
existence of a \(\funp\) proposition.

\paragraph{Conditional model-theoretic results.}
A formally verified theorem may still have explicit premises concerning
Bacon's function domains, invariance, or \QSS.  We therefore distinguish a
verified conditional theorem from a proof that all of its premises hold in
Bacon's particular appendix model.

\section{Goodman's L2 principles}

Goodman's July notes state one principle bearing the label \(L\), in two
forms.  Weak \Ltwo\ is
\[
 \Pure(X)\AND\Pure(Y)\AND\funp(p)\AND\funp(q)\AND Xp=Yq
 \quad\Longrightarrow\quad
 X\approx Y.
\tag{L2}
\]
The strong form requires the same reversible operator to witness both the
relation between the operators and the relation between their inputs:
\[
 \Pure(X)\AND\Pure(Y)\AND\funp(p)\AND\funp(q)\AND Xp=Yq
 \quad\Longrightarrow\quad
 \exists Z\in G\,
   \bigl(X=Y\circ Z\AND q=Zp\bigr).
\tag{strong L2}
\]
The notes contain no separate principle labeled \(L1\).  Accordingly, the
plural ``L principles'' below refers to weak and strong \Ltwo.

\subsection{Object-language verification}

Both forms have been encoded as closed, well-typed object-language formulas.
Isabelle verifies all four uses made in T6:
\[
\begin{array}{c}
 T_0+\PP+\exists\funp+\Ltwo+\Inv\vdash\bot,\\
 T_0+\PP+\exists\funp+\Ltwo+\WI\vdash\bot,\\
 T_0+\PP+\exists\funp+\Ltwo+\TU\vdash\bot,\\
 T_0+\PP+\text{strong \Ltwo}+\RS\vdash\bot.
\end{array}
\]
The repaired Recombination--Exhaustion bridge derives the required
\(\exists\funp\) premise from Recombination, zeroary Exhaustion, and unique
proposition-level fundamentality.  Isabelle also verifies the uses of \Ltwo\
in T7--T9, including uniqueness of kind in T8 and the conditional cardinal
dichotomy in T9.

\subsection{Refutation of L2 in Bacon's appendix model}

Goodman's first suggested attack asks whether \Ltwo\ holds in Bacon's appendix
model when the pure unary operators are exactly the denotations of closed
terms containing only logical vocabulary.  Isabelle answers this question
negatively.

Define the unary operator
\[
 Z(P)=
 \{\,w :
   P(w^\frown 0)\text{ and }P(w^\frown 1)
   \text{ have different truth values}\,\}.
\]
There is a closed term containing only logical vocabulary whose denotation is
exactly \(Z\).  Hence \(Z\) is one of the pure unary operators denoted by
closed terms containing only logical vocabulary.  Isabelle proves:
\[
\begin{array}{ll}
\text{(i)}   & Z\text{ is surjective};\\
\text{(ii)}  & Z(P)=Z(\neg P),\text{ so }Z\text{ is not injective};\\
\text{(iii)} & Z\text{ is right-cancellative on the denotations of closed
                    logical terms};\\
\text{(iv)}  & Z\notin G.
\end{array}
\]
Surjectivity is witnessed explicitly by
\[
 P_S=\{\,v^\frown 0:v\in S\,\},
 \qquad Z(P_S)=S.
\]
Right-cancellativity implies that \(Zp\) is \(\funp\) whenever \(p\) is
\(\funp\).  Choose such a \(p\).  Then
\[
 \mathsf{id}(Zp)=Zp=Z(p),
\]
although \(\mathsf{id}\not\approx Z\), since \(Z\) is not reversible.  This is
an explicit counterexample to \Ltwo.  Since strong \Ltwo\ entails weak
\Ltwo, it fails in this model as well.

\subsection{What this advances}

This settles the first open problem in Goodman's list: \Ltwo\ is false in
Bacon's appendix model under the proposed interpretation of purity.  The
earlier Isabelle analysis had shown
that no counterexample occurs among identity, necessity, possibility,
constant truth, and constant falsity; \(Z\) is the first verified
nonreversible closed logical operator that nevertheless preserves the
distinctions relevant to \(\funp\).

The result does not prove the consistency of \PP.  Bacon's appendix model
still fails \PP\ at the unary-operator type.  What it shows is that failure of
\Ltwo---one of the two escape routes forced by T6---is mathematically
available in Bacon's own semantic setting.  It also shows that a negative
solution cannot treat \Ltwo\ as automatic merely because purity is
interpreted by closed logical definability.

\section{Audit of the object-language theorems}

\renewcommand{\arraystretch}{1.12}
\setlength{\tabcolsep}{4pt}
\begin{longtable}{>{\raggedright\arraybackslash}p{0.55in}
                  >{\raggedright\arraybackslash}p{1.18in}
                  >{\raggedright\arraybackslash}p{4.35in}}
\toprule
\textbf{Item} & \textbf{Status} & \textbf{Result}\\
\midrule
\endfirsthead
\toprule
\textbf{Item} & \textbf{Status} & \textbf{Result}\\
\midrule
\endhead
T1 & Verified &
Under zeroary Exhaustion, every pure proposition is truth or falsity.
Consequently the biconditional operators are identity and negation, and
\WI\ collapses to \Inv.\\

T2a & Verified &
\(\funp\) is closed under every pure reversible operator, hence under
negation.\\

T2b & Verified &
Truth and falsity are not \(\funp\).  A \(\funp\) witness is distinct from
truth, falsity, and its negation.\\

T2c & Verified, stronger &
Every pure proposition is possibly identical to a \(\funp\) witness.  The
proof does not use \PP.\\

T2d & Verified, stronger &
Possible purity of the witness follows without Persistence, although
Goodman's displayed route uses it.\\

T2e & Verified &
The proposition that the witness is noncontingent is false but possible.\\

T2f & Verified conditionally &
The six displayed propositions are pairwise distinct under the explicit
\(\funp(r)\) antecedent.  The antecedent is essential.\\

T3 & Repaired &
The advertised proof has a modal gap: it reaches possible identity where
identity is required.  The theorem follows after adding zeroary Exhaustion,
or from a weaker pure-identity rigidity principle.  A two-world modal
abstraction verifies that the advertised intermediate principles do not by
themselves close the gap.\\

T4 & Verified, stronger &
For each closed pure \(C:t\IC(t\IC t)\),
\(\NOT\funp_{t\IC t}(C(r))\).  The proof uses only the purity schema and
application closure; neither \PP\ nor \(\funp(r)\) is needed.\\

T5 & Verified conditionally &
The assumptions displayed in the notes, together with \(\funp(r)\), yield a \(\funp\) proposition distinct
from \(r\) and \(\NOT r\).  This does not establish consistency of the
antecedent.\\

T6 & Verified &
All four contradiction routes are formalized: \Ltwo+\Inv,
\Ltwo+\WI, \Ltwo+\TU, and strong-\Ltwo+\RS.  Recombination, the
proposition-type instance of Exhaustion, and unique fundamentality derive,
rather than assume, the required \(\funp\) witness where appropriate.  The
advertised T6-\WI\ master claim is also derived directly from its stated
stock, independently of explosion.\\

T7a & Verified &
For the liar family, the identity-indexed member is false, while some
member indexed by a pure reversible operator is true.\\

T7b & Underspecified &
The notes do not give the kind-level diagonal term, its binders, or precise
definitions of ``fixed'' and ``shifted'' point.  There is no unique theorem
to encode.\\

T8 & Verified &
The five base kinds are pairwise distinct; \Ltwo\ gives kind uniqueness; and
the assumptions displayed in the notes prove 31 pairwise-distinct pure unary operators and 31
pairwise-distinct propositions.\\

T9 & Verified at the counting level &
The \PC\ specification makes the map from sets of kinds to pure operators
injective.  Representation of the members of each kind by members of \(G\)
gives an injective code of that kind into \(G\).  These results imply
\[
 |\Pow(K)|\leq |K\times G|.
\]
Classical cardinal arithmetic then gives
\[
 K\text{ finite}\quad\OR\quad |\Pow(K)|\leq |G|.
\]
Instantiating the abstract types uniformly across Goodman's full typed
object language remains separate.\\
\bottomrule
\end{longtable}

\section{Audit of the model-theoretic claims}

\begin{longtable}{>{\raggedright\arraybackslash}p{0.55in}
                  >{\raggedright\arraybackslash}p{1.28in}
                  >{\raggedright\arraybackslash}p{4.25in}}
\toprule
\textbf{Item} & \textbf{Status} & \textbf{Result}\\
\midrule
\endfirsthead
\toprule
\textbf{Item} & \textbf{Status} & \textbf{Result}\\
\midrule
\endhead
M1 & Exact identification; exclusion conditional &
At proposition type, the pure propositions in Bacon's model are exactly the two
invariant propositions, and the operator testing for membership in that class
is the noncontingency operator.  The exact footnote-59 liar is typed, beta-reduced, and proved
pure from \PP\ plus Purity of Fun; given \QSS, this yields a
contradiction.  At the
next type, the infinitary join exists in the full function domain and has
exactly the pure unary operators as its extension.  It interprets
\(\Pure\) at that type.  \PP\ holds there exactly when this operator is
itself in the next pure stock.  The current formalization does not prove
directly that it is absent from the next pure stock in Bacon's appendix
construction.  It proves that absence only from the additional PP-diagonal and
\QSS\ assumptions.  Completeness of the function domain therefore supplies
the join, but does not by itself settle its purity.\\

M2 & Verified &
Goodman's construction \(T\mapsto G_T\) is a bijection between arbitrary sets
of propositions and invariant unary operators.  Cantor's theorem yields more
invariant operators than propositions, so evaluation at one fundamental
proposition cannot be injective on all invariants.  The invariance reading
of purity therefore violates \QSS.\\

M3 & Verified relative to the pure operators &
Relative to the chosen stock of pure operators, \(\funp\) is equivalent to freeness
against nonzero pure laws.  Such propositions have both extreme views.  A
countable Boolean stock has a free generator, while the resulting
\(\funp\) class is topologically meager in the product topology.\\

M4 & Conditional model-theoretic theorem &
Preimage heredity is proved.  The lifted-branch witness is \(\funp\) and
outside the reversible orbit under explicit assumptions: necessitated
\QSS, function-space membership and invariance of every operator in the pure
stock, and presence of identity, zero, and one.  The theorem has not
yet been instantiated for Bacon's stock of operators denoted by closed terms
containing only logical vocabulary.  The
multiple-fundamental ``wide Fun'' argument remains prose-level.\\

M5 & Verified and repaired; one generalization open &
The operator \(f=G_T\) obtained by exchanging \(s\) and \(s'\) is an invariant involution violating \TU, \WI,
and \Inv.  The notes' subsequent countability argument is false: both the
orbit and the displayed world-indexed family are countable.  We replace it
with a uniform diagonal construction.  For every candidate \(R\), it
produces a two-element pair outside every view of \(R\), whose proper views
are all empty or universal.  The induced invariant nonidentity
transposition fixes \(R\).  If it and identity are both included among the
pure operators, \(R\)
ceases to be \(\funp\), proving the corrected pre-rebuild \QSS\
obstruction.  The rebuilding question is now settled.  Isabelle constructs
the least application-closed pure stock containing every closed logical
denotation and the displayed \(f\), proves that \(f\) is a pure self-inverse
operator in the modalized function domain, and rebuilds a fundamental
proposition that supplies Recombination and \(\funp\)-separation at every
world.  The rebuilt interpretation validates Bacon's Recombination
background.  It does not validate \PP; that would require the classifier of
the enlarged pure stock itself to belong to the next pure stock.  The
displayed collision calculation is also verified: the closed operator
\[
  \lambda p\,(p\IFF\mathsf{NC}(p))
\]
takes the same value at \(\mathsf{NC}(r)\) and truth, although
\(\mathsf{NC}(r)\neq\mathsf{truth}\), whenever \(\funp(r)\).  Extending this
argument uniformly to merely existentially given invertibles remains open.\\

M6 & Conditional model-theoretic theorem &
Under the same explicit assumptions---necessitated \QSS, membership of every
pure operator in Bacon's function domain, and invariance---distinct
substitutions are separated by a \(\funp\) proposition.  With identity,
zero, and one, a strict-inclusion \(\funp\) pair blocks the indicated joint
assignment.  Single-coordinate arbitrary realization fails independently
by countability.  These results have not all been instantiated for Bacon's
stock of operators denoted by closed terms containing only logical
vocabulary.\\

M7 & Verified, with an open instance &
The Tarski diagonal lies outside the range of the enumerated closed terms.
For arbitrary invariant operators of Goodman's form \(G_T\), every proposition is reachable from \(r\)
if and only if the orbit map
\[
 i\longmapsto i\cdot r
\]
is injective.  Whether Bacon's glued \(r\) has this property is
construction-sensitive and open.\\

10.1 & Verified, qualified &
Bacon's rebuilt-family theorem is proved in a set-theoretic formalization of
his appendix model.  At type
\(\mathsf{Ind}\), equality is literal identity, and the theorem's hypotheses
therefore force the family to be constant; the theorem cannot supply a family
of distinct individual denotations.\\
\bottomrule
\end{longtable}

\section{Corrections and clarifications to the notes}

\subsection{The modal gap in T3}

The advertised argument uses necessitated \QSS\ and Persistence to move from a
possible fundamental role to an injectivity claim.  The formal derivation
reaches a statement of the form
\[
 \Diamond(Y=Z),
\]
where the desired conclusion requires \(Y=Z\).  No rule of \(H\) licenses
that transition.  The gap can be repaired in either of two principled ways.

\begin{enumerate}
\item Zeroary Exhaustion collapses pure propositions to truth and falsity,
      which supplies the required rigidity.
\item One may assume only the weaker principle that possible identity between
      pure operators implies identity.
\end{enumerate}

\noindent
The unrestricted version of the second principle is false.  Thus the repair
must remain restricted to the pure case.

\subsection{The countability error in M5}

The notes propose choosing a pair \(\{s,s'\}\) outside the orbit of \(r\)
because the orbit is countable and, it is said, there are uncountably many
world-indexed pairs.  But Bacon's worlds are finite sequences of natural
numbers, hence countable.  The displayed family of pairs is therefore
countable as well.  Countability alone proves nothing.

The replacement is direct diagonalization.  Enumerate the views of \(R\) by
the worlds.  Define a proposition \(s_R\) whose truth at the singleton
\([n+2]\) disagrees with the corresponding singleton's membership in the
\(n\)-th view of \(R\), and let
\[
 s'_R=s_R\cup\{\langle\rangle\}.
\]
Membership above a nonactual world depends only on the final symbol of its
sequence.  Hence every nonactual view of \(s_R\) and \(s'_R\) is empty or
universal, while both members differ from every view of \(R\).  Exchanging
their membership in the set \(T_0\) of true propositions yields an invariant
nonidentity operator \(f_R=G_T\) with
\[
 f_R(R)=R.
\]
If \(f_R\) and identity are both included among the pure operators, evaluation at \(R\) is not
injective.  This is exactly the pre-rebuild \QSS\ obstruction the M5 strategy
needs.

The separate rebuilding step is also now verified.  Instead of retaining
this old \(R\), the rebuilt model chooses a new fundamental proposition
generic for the enlarged application-closed stock containing Goodman's
displayed \(f\).  That proposition again satisfies Recombination and
\(\funp\)-separation.

\subsection{The scope of several semantic claims}

The M4 and M6 theorems are genuine mathematical results, but at present they
are conditional.  They assume that every operator counted as pure belongs to
Bacon's function domain and is invariant, together with necessitated \QSS\
and, for some results, the presence of identity and the two constant
operators.
The proofs do not yet instantiate these assumptions for Bacon's complete
stock of operators denoted by closed terms containing only logical
vocabulary.  We therefore do not count them as unconditional
verifications of the corresponding claims about Bacon's particular glued
model.

\section{New theorems advancing the consistency question}

The audit was not merely defensive.  The formalization has produced several
theorems not stated in the original notes.  These results materially clarify
the logical relations among Goodman's central claims.

\subsection{Recombination, Exhaustion, and the existence of a
\texorpdfstring{\(\funp\)}{fun-prime} proposition}

\begin{theorem}
Unary Recombination yields the pointwise \QSS\ core.  Adding zeroary
Exhaustion yields the necessitated universal form needed for \QSS.  Together
with unique proposition-level fundamentality, the resulting stock proves
that a \(\funp\) proposition exists.
\end{theorem}

\noindent
This theorem removes the previously assumed existential \(\funp\) premise
from the premises used in T6.  It also makes explicit exactly where Exhaustion,
rather than Recombination, enters the object-language argument.

\subsection{An exact semantic consequence of T6}

\begin{theorem}
Any model validating Bacon's proof theory, Recombination, the
proposition-type instance of Exhaustion, unique proposition-level
fundamentality, and \PP\ must falsify global \Ltwo\ or global \TU:
\[
 \NOT\,\mathsf{GValid}(\Ltwo)
 \quad\OR\quad
 \NOT\,\mathsf{GValid}(\TU).
\]
Moreover, some world in the constructed model falsifies one of the
two closed formulas.
\end{theorem}

\noindent
Here \(\mathsf{GValid}(A)\) means that the closed formula \(A\) is true at
every world of the model.
This formalizes the central dilemma in Goodman's notes.  It does not select a
disjunct.  A positive model may exploit either cross-kind collisions on
\(\funp\) inputs or a pure reversible operator that is neither uniformly
truth-preserving nor uniformly truth-flipping.

\subsection{The rebuilt M5 model}

\begin{theorem}[Completion of the rebuilding step]
There is a typed interpretation of Bacon's Recombination background whose
pure denotations form the least application-closed stock containing all
closed logical denotations and Goodman's displayed exotic operator \(f\).
That stock is countable.  The operator \(f\) is in the modalized unary
function domain, commutes with taking views, is self-inverse, is not
truth-uniform, and is not a biconditional operator.  A rebuilt fundamental
proposition supplies Recombination and \(\funp\)-separation at every world.
\end{theorem}

\noindent
This discharges the rebuilding step that Goodman's notes left ``modulo
Theorem 10.1.''  It verifies the intended M5 independence result for Bacon's
Recombination background.  It does not provide a model of \PP: \PP\ would
add the further requirement that the classifier of this enlarged pure stock
itself occur at the next type.

\subsection{A uniform pre-rebuild M5 obstruction}

\begin{theorem}[Uniform version of the M5 construction]
For every candidate proposition \(R\), there is a pair
\(\{s_R,s'_R\}\) disjoint from the set of all views of \(R\).  Every
nonactual view of either member is empty or universal.  The associated invariant
transposition fixes \(R\) and is not identity.
\end{theorem}

\noindent
This is stronger and cleaner than the original M5 selection argument.  It
shows that an exotic invariant operator fixing any preassigned \(R\) can be
constructed without cardinal assumptions.  It shows that the pre-rebuild
failure is caused by keeping the old fundamental proposition fixed while
enlarging the pure operators.  The preceding theorem shows how rebuilding
the fundamental proposition avoids that failure.

\subsection{The T9 and M7 reductions}

\begin{theorem}[T9 dichotomy]
If \PC\ supplies an injection \(\Pow(K)\to P\) into the pure unary operators,
and every \(\approx\)-class injects into \(G\), then
\[
 K\text{ is finite}
 \quad\OR\quad
 |G|\geq 2^{|K|}.
\]
\end{theorem}

\begin{theorem}[Invariant reachability]
Every proposition is reachable from \(r\) by an invariant operator if and
only if
\[
 i\longmapsto i\cdot r
\]
is injective.
\end{theorem}

\noindent
The second theorem turns the open invariant form of Fundamental Completeness
into a precise property of the chosen fundamental proposition.  It also
confirms that the issue depends on Bacon's particular construction rather
than following from general facts about substitutions.

\section{Where the open question now stands}

Goodman's question asks whether Bacon's background theory, unique
proposition-level fundamentality, and \PP\ have a model.  The verified results
clarify both possible answers without deciding between them.

\paragraph{Affirmative route.}
A positive answer requires a model of the full theory plus \PP, with purity
interpreted at every relevant type.  T6 shows that such a model must falsify
\Ltwo\ or \TU.  Thus it must contain either pure operators from distinct
\(\approx\)-classes that agree on suitable \(\funp\) propositions, or a pure
reversible operator that is neither uniformly truth-preserving nor uniformly
truth-flipping.  Bacon's substitution-structure models do not already supply
the desired example: M1 locates their first obstruction to \PP, while M2 shows
that identifying purity with invariance is incompatible with \QSS.  The
operator \(Z\) shows that the operators denoted by Bacon's closed logical
terms already take the \(\neg\Ltwo\) branch, but the model's independent
failure of \PP\ prevents
this from being a positive solution.

\paragraph{Negative route.}
The most direct negative route is the one already isolated in Goodman's
notes: derive enough of \Ltwo\ and one of \Inv, \WI, \TU, or \RS\ from the
remaining principles, and apply T6.  The formalization verifies the relevant
T6 implications and removes the need to assume a \(\funp\) witness separately.
The \(Z\) counterexample shows that Bacon's base theory and the interpretation
of purity by denotations of closed logical terms do not by themselves secure
\Ltwo.  Any
negative proof along this route must therefore use \PP\ or some further
principle not validated by that model.

\paragraph{Relative plausibility.}
The formal work does not establish that either answer is more likely.  It
shows exactly what each answer must overcome.  A positive model must depart
from Bacon's familiar substitution structures in one of the ways forced by
the T6 disjunction.  A negative proof must promote one of Goodman's
model-sensitive principles, especially \Ltwo\ or \TU, to a consequence of the
background theory itself.

\section{Remaining verification work on the notes}

The following items should not be reported as verified:

\begin{enumerate}
\item T7b, until a precise kind-level diagonal formula is supplied.
\item The multiple-fundamental ``wide Fun'' discussion in M4, which does not
      state a precise selection condition on \(\Fun\) or a precise Separated
      Structure formula.
\item The proposed extension of the verified M5 collision calculation from
      its displayed operator to merely existentially given invertibles.
\end{enumerate}

\section{Recommended next steps}

The following tasks arise directly from the unresolved or qualified claims in
Goodman's notes.

\begin{enumerate}
\item \textbf{Obtain determinate formulations.}
      Obtain an exact T7b statement, an exact formulation of the wide-Fun
      discussion before encoding them.
\item \textbf{Generalize the M5 collision method.}
      Determine whether the displayed noninjectivity calculation can be
      extended to invertible operators given only existentially, as would be
      required for the proposed PP-derived use of the method.
\item \textbf{Test Goodman's negative route.}
      In light of the \(Z\) counterexample, determine whether \Ltwo\ follows
      only after adding \PP\ or some independently acceptable principle not
      valid in Bacon's appendix model.  Together with \TU, \WI, \Inv, or the
      strong-\Ltwo+\RS\ route, the verified T6 theorems then give the
      corresponding inconsistency result.
\item \textbf{Keep model claims and proof-theoretic claims separate.}
      In particular, do not promote the conditional M4/M6 stock theorems to
      claims about Bacon's appendix model until their stock premises are
      instantiated.
\end{enumerate}

\section{Reproducibility}

The repository's three-layer organization and current build commands are
documented in \isa{README.md} and
\isa{docs/REPOSITORY_STRUCTURE.md}.  We do not reproduce that implementation
documentation here.

The controlling verification ledger is
\[
\isa{reports/GOODMAN_COMPLETE_VERIFICATION_MATRIX_2026-07-27.md}.
\]
The dedicated audit theory is
\[
\isa{reports/audit_goodman_complete/Audit_Goodman_Complete.thy}.
\]
The verification matrix gives, for each claim, the exact Isabelle theorem and
the qualification under which it has been proved.  The project status and
handoff files identify the source theory corresponding to each result.

\bigskip

\noindent
The appropriate overall verdict is therefore:
\emph{the formalization verifies, corrects, or precisely qualifies every
determinate claim in Goodman's notes, including two substantive repairs.  It
also completes the rebuilding step in M5, refutes L2 in Bacon's appendix
model under the proposed interpretation of purity, and proves the exact T6
semantic consequence and several strengthening results directly relevant to
the notes.  The remaining underspecified and model-dependent claims are
explicitly identified.  Goodman's consistency question itself remains open.}

\end{document}
